\section{Exact arithmetic certificate and its soundness}\label{app:arithmetic}

This appendix specifies all numerical operations needed in the proof and
proves that their outputs enclose the corresponding real expressions.
The numerical part of the argument is a finite sequence of integer
inequalities; it does not assume correctness of a floating-point error
estimate. The source files and complete rational records accompany the paper.

\subsection{Dyadic representation and elementary operations}
Fix $P=160$ and $D=2^P$. An enclosure is a pair of integers $(a,b)$,
$a\le b$, representing $[a/D,b/D]$.
An exact fraction $s/t$, $t>0$, is enclosed by
\[
 \left(\left\lfloor\frac{sD}{t}\right\rfloor,
       \left\lceil\frac{sD}{t}\right\rceil\right).
\]
Addition and negation act exactly on endpoints. To multiply two enclosures,
take the minimum and maximum of the four integer endpoint products, divide
by $D$, and round the lower endpoint down and the upper endpoint up.
If $0<a\le b$, the reciprocal enclosure is
\[
 \left(\left\lfloor\frac{D^2}{b}\right\rfloor,
       \left\lceil\frac{D^2}{a}\right\rceil\right).
\]
For a strictly negative interval use sign reversal; division is rejected if
the denominator interval contains zero. Integer powers are evaluated by
repeated squaring, using these same enclosure operations.
For square roots of $0\le a\le b$, let
\[
 u=\lfloor\sqrt{aD}\rfloor,\qquad v=\lfloor\sqrt{bD}\rfloor.
\]
The integer witnesses $u^2\le aD<(u+1)^2$ and
$v^2\le bD<(v+1)^2$ are checked explicitly. The resulting enclosure is
$(u,v)$ if $v^2=bD$, and $(u,v+1)$ otherwise.

For intervals $X_i=[a_i,b_i]$, the minimum and maximum of their unknown
values are enclosed, respectively, by
$[\min_i a_i,\min_i b_i]$ and $[\max_i a_i,\max_i b_i]$.
These rules do not require independence among the unknown values.
Likewise, reusing an uncertain expression may widen an interval but cannot
invalidate its enclosure. A strict comparison $X<Y$ is accepted only when
the upper endpoint of $X$ is below the lower endpoint of $Y$.

\begin{lemma}[Elementary enclosure invariant]\label{lem:invariant}
Every elementary operation just specified encloses its real result whenever
its operands are valid enclosures and its stated domain conditions hold.
\end{lemma}
\begin{proof}
For multiplication, a bilinear function attains its extrema on a rectangle
at corners. Reciprocal is decreasing on each of its sign domains, square
root is increasing on $[0,\infty)$, and minimum and maximum are coordinatewise
increasing. The floor and ceiling definitions preserve the appropriate
one-sided bounds. Addition and negation are immediate. Induction on the
number of multiplications proves the integer-power claim.
\end{proof}

\subsection{Exponential enclosures}
For a nonnegative exact rational $x$, choose an integer $r\ge0$ so that
$y=x/2^r\le1/8$. With $N=40$, define the exact fractions
\[
 S_N(y)=\sum_{k=0}^N\frac{y^k}{k!},\qquad
 R_N(y)=\frac{y^{N+1}}{(N+1)!}\frac1{1-y/(N+2)}.
\]
Then
\begin{equation}\label{eq:exp-enclosure}
 S_N(y)\le e^y\le S_N(y)+R_N(y).
\end{equation}
Indeed, all terms are nonnegative, and the ratio of consecutive terms after
the first omitted one is at most $y/(N+2)<1$, so a geometric series bounds
the tail. Enclose both endpoints as dyadic fractions and square the result
$r$ times. This encloses $e^x$ by Lemma~\ref{lem:invariant}. For negative
$x$, enclose $e^{-x}$ first and take the reciprocal. For an input interval
$[a,b]$, enclose $e^a$ and $e^b$ separately and retain the outer endpoints,
using monotonicity of $\exp$.

All Taylor coefficients and remainders in the implementation are computed
with exact fractions. Thus the finite series length affects enclosure width,
not validity. Very small positive exponentials may be enclosed by an interval
whose lower endpoint is zero; this is a valid enclosure, not a replacement
of a positive real number by zero. A subsequent division by such an interval
would be rejected.

\subsection{Logarithm and circle-constant enclosures}
Let $z=1/3$ and $N=80$. Integrating the geometric series for $1/(1-z^2)$
gives
\[
 \log2=2\sum_{k=0}^\infty\frac{z^{2k+1}}{2k+1}.
\]
Its remainder after $N$ terms is positive and at most
\begin{equation}\label{eq:log-enclosure}
 \frac{2z^{2N+1}}{(2N+1)(1-z^2)}.
\end{equation}
Consequently this partial sum and remainder give an exact rational enclosure
for $\log2$. Squaring and multiplying by 144 gives one for $C_*$.
At the critical scale the identity $2t_0=1$ is simplified symbolically
before evaluation: for $C=C_*$,
\[
 t_j=\frac{(j+3)^2}{18\cdot2^j}.
\]
Thus the endpoint sign condition does not depend on rounding a number close
to one.

To enclose $\pi$ we use Machin's identity
\begin{equation}\label{eq:machin}
 \pi=16\arctan(1/5)-4\arctan(1/239).
\end{equation}
Here is a direct verification. If $\theta=\arctan(1/5)$ and
$\phi=\arctan(1/239)$, double-angle formulas give
$\tan(2\theta)=5/12$ and $\tan(4\theta)=120/119$; therefore
$\tan(4\theta-\phi)=1$. Also $0<4\theta-\phi<4/5<\pi/2$:
positivity follows from $\theta>\phi>0$, and
$\theta<1/5$ follows from the integral formula for arctangent.
The elementary bound $\pi/2>1$ follows from
$\pi/4=\int_0^1(1+x^2)^{-1}dx>1/2$.
Injectivity of tangent on this interval gives $4\theta-\phi=\pi/4$.

For $z\in\{1/5,1/239\}$, the alternating expansion
\[
 \arctan z=\sum_{k=0}^{N-1}(-1)^k\frac{z^{2k+1}}{2k+1}+R_N,
 \qquad 0\le R_N\le\frac{z^{2N+1}}{2N+1},\qquad N=50,
\]
follows by integrating the finite geometric expansion of $1/(1+x^2)$;
$N$ is even. Combining these rational enclosures in \eqref{eq:machin}
encloses $\pi$.

\begin{lemma}[Soundness of the certificate arithmetic]\label{lem:soundness}
Every enclosure produced by \texttt{rational\_interval.py} for an expression
used in the certificate contains the exact real value of that expression.
Every strict numerical inequality accepted by \texttt{exact\_certificate.py}
is therefore a true real inequality.
\end{lemma}
\begin{proof}
The leaf constants are exact rationals, the logarithm enclosure
\eqref{eq:log-enclosure}, and the circle-constant enclosure
\eqref{eq:machin}. They are sound by the preceding arguments.
Structural induction on the expression tree uses Lemma~\ref{lem:invariant}
for elementary operations and \eqref{eq:exp-enclosure} for exponentials.
Every division checks separation from zero, every square root checks its
nonnegative domain and integer witnesses, and conditional use of
$1/(1-\Theta_j)$ requires a certified strict inequality $\Theta_j<1$.
The last assertion follows from disjoint interval endpoints.
\end{proof}

\subsection{Finite specification of the obligations}
The verification comprises the following explicitly bounded operations.
\begin{enumerate}
\item Establish $693/1000<\log2<694/1000$, $69<C_*<70$, the parameter
inequalities, and the four positive Taylor inequalities in the counting proof.
\item For each of the two values $C\in\{C_*,70\}$ and each integer
$i\in\{150,\ldots,319\}$, evaluate the slab expressions on
$[i/10,(i+1)/10]$. Include scales until the first $j+1$ for which
$2^{j+1}$ is certified greater than $256((i+1)/10)^2/C$.
This can include an unnecessary scale if a comparison is inconclusive;
including additional nonnegative error terms is conservative.
\item For each $C$, evaluate the eight low-scale expressions at $H=32$,
check their derivative conditions, and evaluate the omitted-scale bound.
For $C=C_*$ use the exact rational formula for $t_j$ above.
\item For each row of Table~\ref{tab:constants}, take maxima over all slabs
starting at or above its $L_0$, the corresponding half-line bound, and
$1/1000$. Check all three required error bounds against the printed rational
$B$. Check that $B+10^{-9}$ is still below the failure threshold associated
with the table's convenient success guarantee.
\item Compute each amplification probability as an exact finite sum of
fractions, and compare it with $1/100$.
\end{enumerate}
All loops have explicit finite limits, except scale inclusion, which terminates
because its fixed finite comparison bound is eventually exceeded by $2^j$.
The exact intervals are retained in the certificate file. In particular,
there is no extrapolation from a sampled distribution, no test of only
selected integer $n$, and no truncation of the half-line $L\ge32$.

\subsection{Independent risk aggregation and reproducibility}
For every slab the certificate records rational upper bounds
$v_j^+,w_j^+,h^+,o^+$ for $v_j,w_j,H(L),2/n$.
The separate script \texttt{audit\_exact.py} reconstructs, using the standard
exact Fraction type,
\begin{align*}
 U^+&=\sum_j\frac{v_j^+}{v_j^++d_j^2},\\
 R_f^+&=\max\left\{h^+,\max_j\frac{v_j^+}{v_j^++\kappa_j^2}\right\},\\
 R^+&=\max\left\{h^+,o^+,\max_j\frac{2w_j^+}{2w_j^++\kappa_j^2}\right\},\\
 B^+&=\max\{U^++R_f^+,R^+\}.
\end{align*}
These are valid bounds because the displayed rational functions are
coordinatewise nondecreasing for nonnegative arguments.
The auditor similarly reconstructs the half-line aggregation, checks that
the complete set of slabs is present, verifies the source hashes, and checks
all claimed thresholds. It also tests the arithmetic operations against
independent Fraction calculations and checks exact Chebyshev coefficient
identities. Those tests complement the mathematical soundness proof;
they do not replace it.

From the artifact directory, the commands are
\begin{verbatim}
python3 exact_certificate.py
python3 audit_exact.py
tectonic submission_proof.tex
\end{verbatim}
Only Python's standard library is needed for the first two commands.
The checker rejects float inputs and does not call a logarithm, exponential,
or trigonometric routine from a numerical library. Its checks raise explicit
exceptions and remain active under Python's optimization flag.

\paragraph{Scope of formalization.}
The theorem, reduction to a finite family of rational inequalities, and
soundness of the enclosure rules are mathematical proofs supplied in this
document. The artifact executes those finite checks. There is no claim that
the complete probability and real-analysis development has been encoded in
Lean, Coq, Isabelle, or another foundational proof assistant, and no new
axiom asserting the desired statistical theorem is introduced. A full
proof-assistant formalization would be an additional deliverable.
